
\documentclass{article}
\usepackage{colortbl}
\usepackage{makecell}
\usepackage{multirow}
\usepackage{supertabular}

\begin{document}

\newcounter{utterance}

\centering \large Interaction Transcript for game `cladder', experiment `full\_v1.5\_default', episode 9763 with refineonly{-}llama{-}3.1{-}8b{-}instruct.
\vspace{24pt}

{ \footnotesize  \setcounter{utterance}{1}
\setlength{\tabcolsep}{0pt}
\begin{supertabular}{c@{$\;$}|p{.15\linewidth}@{}p{.15\linewidth}p{.15\linewidth}p{.15\linewidth}p{.15\linewidth}p{.15\linewidth}}
    \# & \multicolumn{2}{c}{Player} && \multicolumn{2}{c}{Game Master} \\
    \hline

    \theutterance \stepcounter{utterance}  
    & & & \multicolumn{4}{p{0.6\linewidth}}{
        \cellcolor[rgb]{0.9,0.9,0.9}{
            \makecell[{{p{\linewidth}}}]{
                \texttt{\tiny{[P1$\langle$GM]}}
                \texttt{You are an expert in causal inference. The following question is not a typical commonsense query, but rather a meticulously designed question created by a professor specializing in causal inference, intended to assess the students' mastery of the course content.} \\
\texttt{Imagine a self{-}contained, hypothetical world with only the following conditions, and without any unmentioned factors or causal relationships: Yupt has a direct effect on jyka and kwox. Jyka has a direct effect on kwox. We know that yupt causes jyka. yupt and jyka causes kwox. We observed an individual is yupt. Would an individual is kwox if not jyka instead of jyka?} \\
\texttt{Start your answer with 'Yes' or 'No', followed by additional reasoning or evidence to support your explanation.} \\
            }
        }
    }
     \\ \\

    \theutterance \stepcounter{utterance}  
    & \multicolumn{4}{p{0.6\linewidth}}{
        \cellcolor[rgb]{0.561,0.737,0.561}{
            \makecell[{{p{\linewidth}}}]{
                \texttt{\tiny{[P1$\rangle$GM]}}
                \texttt{No.} \\
\\ 
\texttt{The given conditions imply a causal graph with the following structure:} \\
\\ 
\texttt{{-} Yupt {-}> Jyka} \\
\texttt{{-} Yupt {-}> Kwox} \\
\texttt{{-} Jyka {-}> Kwox} \\
\\ 
\texttt{We are asked to compare the counterfactual outcome of Kwox if Jyka were not present instead of being present. This is a classic problem of counterfactual inference.} \\
\\ 
\texttt{To solve this, we can use the backdoor criterion, which states that if we can block all backdoor paths between the treatment (Jyka) and the outcome (Kwox) using observed variables, then we can estimate the causal effect of Jyka on Kwox.} \\
\\ 
\texttt{In this case, the backdoor path is Yupt {-}> Jyka {-}> Kwox. We can block this path by conditioning on Yupt, which is the common cause of both Jyka and Kwox.} \\
\\ 
\texttt{However, the presence of Yupt does not allow us to block the backdoor path because Yupt is a direct cause of Kwox. This means that the causal effect of Jyka on Kwox is not identifiable using the backdoor criterion.} \\
\\ 
\texttt{To further analyze this problem, we can use the concept of causal mediation. The effect of Yupt on Kwox is mediated by both Jyka and the direct effect of Yupt on Kwox. If we were to remove Jyka, the effect of Yupt on Kwox would be reduced, but not eliminated, because Yupt still has a direct effect on Kwox.} \\
\\ 
\texttt{Therefore, the counterfactual outcome of Kwox if Jyka were not present instead of being present is not necessarily zero. However, the exact value of this counterfactual outcome is not identifiable from the given information.} \\
\\ 
\texttt{In summary, the answer is No, but the reasoning is more nuanced than a simple yes or no. The causal effect of Jyka on Kwox is not identifiable using the backdoor criterion, and the counterfactual outcome of Kwox if Jyka were not present is not necessarily zero.} \\
            }
        }
    }
    & & \\ \\

    \theutterance \stepcounter{utterance}  
    & & & \multicolumn{2}{p{0.3\linewidth}}{
        \cellcolor[rgb]{0.9,0.9,0.9}{
            \makecell[{{p{\linewidth}}}]{
                \texttt{\tiny{[GM$|$GM]}}
                \texttt{No} \\
            }
        }
    }
    & & \\ \\

    \theutterance \stepcounter{utterance}  
    & & & \multicolumn{2}{p{0.3\linewidth}}{
        \cellcolor[rgb]{0.9,0.9,0.9}{
            \makecell[{{p{\linewidth}}}]{
                \texttt{\tiny{[GM$|$GM]}}
                \texttt{no} \\
            }
        }
    }
    & & \\ \\

    \theutterance \stepcounter{utterance}  
    & & & \multicolumn{2}{p{0.3\linewidth}}{
        \cellcolor[rgb]{0.9,0.9,0.9}{
            \makecell[{{p{\linewidth}}}]{
                \texttt{\tiny{[GM$|$GM]}}
                \texttt{game\_result = WIN} \\
            }
        }
    }
    & & \\ \\

\end{supertabular}
}

\end{document}
